\documentclass[aspectratio=169,12pt]{beamer}
% Estilo canónico para las presentaciones de Métodos Numéricos 2026-2.
\usepackage[T1]{fontenc}
\usepackage[utf8]{inputenc}
\usepackage[spanish,es-nodecimaldot]{babel}
\usepackage{lmodern}
\usepackage{amsmath,amssymb,mathtools}
\usepackage{booktabs}
\usepackage{graphicx}
\usepackage{xcolor}
\usepackage{hyperref}

\definecolor{upazul}{HTML}{003B70}
\definecolor{upazulclaro}{HTML}{EAF2F8}
\definecolor{updorado}{HTML}{C5962E}
\definecolor{upgris}{HTML}{4A5560}
\definecolor{upverde}{HTML}{167C5A}
\definecolor{uprojo}{HTML}{B33A3A}

\setbeamercolor{normal text}{fg=black,bg=white}
\setbeamercolor{structure}{fg=upazul}
\setbeamercolor{frametitle}{fg=white,bg=upazul}
\setbeamercolor{title}{fg=upazul}
\setbeamercolor{subtitle}{fg=upgris}
\setbeamercolor{block title}{fg=white,bg=upazul}
\setbeamercolor{block body}{fg=black,bg=upazulclaro}
\setbeamercolor{alerted text}{fg=uprojo}
\setbeamercolor{example text}{fg=upverde}

\setbeamerfont{title}{size=\Huge,series=\bfseries}
\setbeamerfont{subtitle}{size=\Large}
\setbeamerfont{frametitle}{size=\LARGE,series=\bfseries}
\setbeamerfont{normal text}{size=\large}
\setbeamerfont{block title}{series=\bfseries}

\setbeamertemplate{navigation symbols}{}
\setbeamertemplate{itemize item}{\color{updorado}$\blacktriangleright$}
\setbeamertemplate{itemize subitem}{\color{upazul}$\bullet$}
\setbeamertemplate{blocks}[rounded][shadow=false]

\setbeamertemplate{frametitle}{%
  \nointerlineskip
  \begin{beamercolorbox}[wd=\paperwidth,ht=0.78cm,dp=0.24cm,leftskip=0.55cm]{frametitle}
    \usebeamerfont{frametitle}\insertframetitle
  \end{beamercolorbox}}

\setbeamertemplate{footline}{%
  \leavevmode
  \hbox{%
    \begin{beamercolorbox}[wd=.82\paperwidth,ht=2.7ex,dp=1.2ex,leftskip=.5cm]{author in head/foot}
      \color{upgris}\scriptsize Métodos Numéricos 2026-2 \hspace{0.8em} | \hspace{0.8em} Universidad del Pacífico
    \end{beamercolorbox}%
    \begin{beamercolorbox}[wd=.18\paperwidth,ht=2.7ex,dp=1.2ex,rightskip=.5cm plus1fil]{date in head/foot}
      \color{upazul}\scriptsize\hfill\insertframenumber/\inserttotalframenumber
    \end{beamercolorbox}}
  \vskip0pt}

\newcommand{\concepto}[1]{\textcolor{upazul}{\textbf{#1}}}
\newcommand{\resultado}[1]{\textcolor{upverde}{\textbf{#1}}}
\newcommand{\advertencia}[1]{\textcolor{uprojo}{\textbf{#1}}}
\newcommand{\codigo}[1]{\colorbox{upazulclaro}{\strut\texttt{#1}}}
\newcommand{\fuente}[1]{\vfill\hfill{\tiny\color{upgris}#1}}


\title{Aproximación lineal}
\subtitle{Mínimos cuadrados y ecuaciones normales}
\author{Oswaldo José Velásquez Castañón}
\institute{Universidad del Pacífico}
\date{Métodos Numéricos -- 2026-2}
\newcommand{\R}{\mathbb{R}}
\newcommand{\mat}[1]{\begin{bmatrix}#1\end{bmatrix}}
\newcommand{\norm}[1]{\left\lVert#1\right\rVert}
\begin{document}
\begin{frame}[plain]
\begin{columns}[T,onlytextwidth]
\begin{column}{.70\textwidth}
\vspace{.7cm}{\usebeamerfont{title}\color{upazul}\inserttitle\par}
\vspace{.35cm}{\usebeamerfont{subtitle}\color{upgris}\insertsubtitle\par}
\vspace{1cm}{\large\insertauthor\par}{\normalsize\insertdate\par}
\end{column}
\begin{column}{.25\textwidth}\vspace{.55cm}\includegraphics[width=\linewidth]{UPlogo.pdf}\end{column}
\end{columns}
\end{frame}

\begin{frame}{Objetivos de la clase}
Al terminar esta clase podremos:
\begin{itemize}
\item construir una matriz de diseño a partir de datos y funciones base;
\item formular el ajuste como un problema de mínimos cuadrados;
\item deducir las ecuaciones normales y justificar su suficiencia;
\item interpretar el residuo como una componente ortogonal;
\item distinguir unicidad de los coeficientes y unicidad del ajuste.
\end{itemize}
\end{frame}

\begin{frame}{Datos y modelo}
Supongamos que observamos cantidades $x_i$ y costos $y_i$:
\begin{center}\begin{tabular}{ccccc}\toprule
$x_i$ & 0 & 1 & 2 & 3\\
$y_i$ & 1 & 2 & 2 & 4\\\bottomrule
\end{tabular}\end{center}
Buscamos un modelo sencillo
\[
\widehat y(x)=c+ax.
\]
Los datos de esta clase son ilustrativos. El parámetro $c$ representa el costo
estimado cuando $x=0$ y $a$ el incremento por unidad de $x$.
\end{frame}

\begin{frame}{Interpolación y ajuste}
\begin{center}\begin{tabular}{p{5.5cm}p{5.5cm}}\toprule
Interpolación & Mínimos cuadrados\\\midrule
Exige $p(x_i)=y_i$ en todos los nodos. & Elige un modelo que minimiza un error agregado.\\[2mm]
La cantidad de parámetros permite cumplir las condiciones. & Puede haber muchas más observaciones que parámetros.\\[2mm]
Es útil para representar valores de una función conocida. & Es útil cuando hay ruido o se busca una descripción sencilla.\\\bottomrule
\end{tabular}\end{center}
Una recta tiene dos parámetros. En general, no puede atravesar cuatro puntos arbitrarios.
\end{frame}

\begin{frame}{Residuos y función objetivo}
El residuo de la observación $i$ es
\[
r_i=y_i-(c+ax_i).
\]
Minimizamos la suma de cuadrados de los residuos:
\[
S(c,a)=\sum_{i=1}^m\bigl(y_i-c-ax_i\bigr)^2.
\]
\begin{block}{Qué distancia medimos}
Esta formulación usa desviaciones verticales en $y$. Las distancias
perpendiculares a la recta corresponden a otro problema.
\end{block}
\end{frame}

\begin{frame}{Formulación matricial}
Definimos
\[
A=\mat{1&x_1\\\vdots&\vdots\\1&x_m},
\quad \theta=\mat{c\\a},\quad b=\mat{y_1\\\vdots\\y_m}.
\]
Entonces $r=b-A\theta$ y el problema es
\[
\boxed{\min_{\theta\in\R^2}\norm{A\theta-b}_2^2.}
\]
La matriz $A$ contiene los valores de las funciones base $1$ y $x$.
\end{frame}

\begin{frame}{Linealidad en los coeficientes}
La aproximación general tiene la forma
\[
\widehat y(x)=\sum_{j=1}^n\theta_j\varphi_j(x),
\qquad A_{ij}=\varphi_j(x_i).
\]
Las funciones $\varphi_j$ pueden ser no lineales en $x$:
\[
1,\ x,\ x^2,\ \sin x,\ \cos x,\ e^x.
\]
\begin{block}{Aproximación lineal}
La combinación es lineal en los coeficientes desconocidos $\theta_j$.
Un ajuste polinómico también es un problema de mínimos cuadrados lineales.
\end{block}
\end{frame}

\begin{frame}{Sistemas sobredeterminados}
Para $A\in\R^{m\times n}$, con $m\geq n$, consideramos
\[
\min_{\theta\in\R^n}\norm{A\theta-b}_2.
\]
Si $m>n$, hay más ecuaciones que incógnitas. El sistema $A\theta=b$ puede
ser incompatible.
\begin{block}{Solución aproximada}
Buscamos el vector $A\theta$ más cercano a $b$ entre todos los vectores del
espacio columna de $A$.
\end{block}
Minimizar la norma o su cuadrado da los mismos minimizadores.
\end{frame}

\begin{frame}{Deducción de las ecuaciones normales}
Definimos
\[
J(\theta)=\frac12\norm{A\theta-b}_2^2
=\frac12\theta^TA^TA\theta-\theta^TA^Tb+\frac12b^Tb.
\]
Su gradiente y su hessiano son
\[
\nabla J(\theta)=A^T(A\theta-b),\qquad
\nabla^2J(\theta)=A^TA.
\]
Una solución satisface
\[
\boxed{A^TA\theta^*=A^Tb.}
\]
\end{frame}

\begin{frame}{Teorema de aproximación lineal}
Sea $A\in\R^{m\times n}$, $m\geq n$, y $b\in\R^m$. Entonces:
\begin{enumerate}
\item $\theta^*$ minimiza $\norm{A\theta-b}_2$ si y solo si
$A^TA\theta^*=A^Tb$.
\item El conjunto de minimizadores es afín y no vacío. Todos sus elementos
producen el mismo vector $A\theta^*$.
\item El minimizador es único si y solo si $\operatorname{rango}(A)=n$.
\item Entre los minimizadores existe exactamente uno de norma euclídea mínima.
\end{enumerate}
\end{frame}

\begin{frame}{La condición también es suficiente}
Si $A^T(A\theta^*-b)=0$, para cualquier perturbación $h$,
\[
\begin{aligned}
J(\theta^*+h)
&=J(\theta^*)+h^TA^T(A\theta^*-b)+\tfrac12\norm{Ah}_2^2\\
&=J(\theta^*)+\tfrac12\norm{Ah}_2^2\geq J(\theta^*).
\end{aligned}
\]
\begin{block}{Teorema}
$\theta^*$ minimiza $\norm{A\theta-b}_2$ si y solo si satisface las ecuaciones normales.
\end{block}
La convexidad del problema permite pasar de estacionariedad a mínimo global.
\end{frame}

\begin{frame}{Geometría del residuo}
Con $r^*=b-A\theta^*$, las ecuaciones normales equivalen a
\[
A^Tr^*=0.
\]
Por tanto, el residuo es ortogonal a cada columna de $A$.
\[
b=\underbrace{A\theta^*}_{\text{proyección en }\mathcal C(A)}
+\underbrace{r^*}_{\text{componente ortogonal}}.
\]
La identidad de Pitágoras muestra que cualquier cambio dentro de $\mathcal C(A)$
aumenta o mantiene la distancia a $b$.
\end{frame}

\begin{frame}{Rango y unicidad}
Para cualquier $v\in\R^n$,
\[
v^TA^TAv=\norm{Av}_2^2.
\]
En consecuencia,
\[
\mathcal N(A^TA)=\mathcal N(A).
\]
\begin{block}{Coeficientes únicos}
El minimizador es único si y solo si $\operatorname{rango}(A)=n$.
En ese caso $A^TA$ es definida positiva.
\end{block}
Con rango deficiente, varias combinaciones de coeficientes producen el mismo ajuste.
\end{frame}

\begin{frame}{Otras normas de aproximación}
Con norma $1$, introducimos $t\in\R^m$:
\[
\min_{\theta,t}\sum_i t_i,\qquad
-t\leq A\theta-b\leq t,\quad t\geq0.
\]
Con norma infinito, introducimos un escalar $t$:
\[
\min_{\theta,t}t,\qquad
-t\mathbf1\leq A\theta-b\leq t\mathbf1,\quad t\geq0.
\]
Ambas formulaciones son problemas de programación lineal.
La norma $2$ conduce a las ecuaciones normales estudiadas en esta clase.
\end{frame}

\begin{frame}{La recta de mínimos cuadrados}
Para $\widehat y=c+ax$, las ecuaciones normales son
\[
\mat{m&\sum_i x_i\\\sum_i x_i&\sum_i x_i^2}
\mat{c\\a}=\mat{\sum_i y_i\\\sum_i x_iy_i}.
\]
Si los $x_i$ no son todos iguales,
\[
a=\frac{\sum_i(x_i-\bar x)(y_i-\bar y)}{\sum_i(x_i-\bar x)^2},
\qquad c=\bar y-a\bar x.
\]
La recta ajustada pasa por $(\bar x,\bar y)$.
\end{frame}

\begin{frame}{Ejemplo: construcción del sistema}
Con los cuatro datos de la clase,
\[
A=\mat{1&0\\1&1\\1&2\\1&3},\qquad b=\mat{1\\2\\2\\4}.
\]
Las sumas necesarias son
\[
\sum_i x_i=6,\quad\sum_i x_i^2=14,\quad
\sum_i y_i=9,\quad\sum_i x_iy_i=18.
\]
Por tanto,
\[
\mat{4&6\\6&14}\mat{c\\a}=\mat{9\\18}.
\]
\end{frame}

\begin{frame}{Ejemplo: solución y residuos}
Al resolver las ecuaciones normales obtenemos $c=0.9$ y $a=0.9$.
\begin{center}\begin{tabular}{rrrrr}\toprule
$x_i$ & $y_i$ & $\widehat y_i$ & $r_i$ & $r_i^2$\\\midrule
0 & 1 & 0.9 & 0.1 & 0.01\\
1 & 2 & 1.8 & 0.2 & 0.04\\
2 & 2 & 2.7 & $-0.7$ & 0.49\\
3 & 4 & 3.6 & 0.4 & 0.16\\\bottomrule
\end{tabular}\end{center}
\[
S=\sum_i r_i^2=0.70,\qquad \norm{r}_2=\sqrt{0.70}\approx0.836660.
\]
\end{frame}

\begin{frame}{Verificación de la solución}
Para $r=(0.1,0.2,-0.7,0.4)^T$,
\[
\sum_i r_i=0,\qquad \sum_i x_ir_i=0.
\]
Esto verifica $A^Tr=0$.
\begin{block}{Dos propiedades de una recta con intercepto}
La suma de residuos es cero y el residuo es ortogonal al predictor.
Estas propiedades caracterizan el óptimo en aritmética exacta.
\end{block}
Un residuo no nulo es compatible con una solución correcta de mínimos cuadrados.
\end{frame}

\begin{frame}{Existencia con rango deficiente}
La proyección de $b$ sobre $\mathcal C(A)$ existe siempre y es única.
Por eso siempre hay al menos un minimizador.
\[
\{\theta:\theta\text{ minimiza }J\}
=\theta^*+\mathcal N(A).
\]
Si $z\in\mathcal N(A)$, entonces
\[
A(\theta^*+z)=A\theta^*.
\]
Los coeficientes pueden variar, mientras los valores ajustados y el residuo
permanecen iguales.
\end{frame}

\begin{frame}{Existencia: prueba algebraica}
Un sistema $M\theta=d$ es compatible si y solo si
\[
z^Td=0\qquad\text{para todo }z\in\mathcal N(M^T).
\]
Para $M=A^TA$ y $d=A^Tb$, sea $z\in\mathcal N(A^TA)$.
Entonces
\[
0=z^TA^TAz=\norm{Az}_2^2,\qquad Az=0.
\]
Por tanto,
\[
z^TA^Tb=(Az)^Tb=0.
\]
Las ecuaciones normales son compatibles para cualquier $b$.
\end{frame}

\begin{frame}{Conjunto afín y ajuste único}
Sea $\theta_0$ una solución de las ecuaciones normales. Entonces
\[
\mathcal L=\{\theta:A^TA\theta=A^Tb\}
=\theta_0+\mathcal N(A^TA)=\theta_0+\mathcal N(A).
\]
Si $\theta_1,\theta_2\in\mathcal L$, su diferencia pertenece a $\mathcal N(A)$.
Por ello,
\[
A\theta_1=A\theta_2.
\]
Si el núcleo es trivial, $\mathcal L$ contiene un solo elemento.
Si el núcleo contiene un vector no nulo, hay infinitas soluciones.
\end{frame}

\begin{frame}{Existencia de norma mínima}
El conjunto $\mathcal L$ es cerrado, afín y no vacío. Elegimos $\theta_0\in\mathcal L$.
El conjunto
\[
K=\mathcal L\cap\{\theta:\norm{\theta}_2\leq\norm{\theta_0}_2\}
\]
es compacto y no vacío. La función continua $\theta\mapsto\norm{\theta}_2$
alcanza un mínimo en $K$.
\begin{block}{Mínimo en todo $\mathcal L$}
Los elementos de $\mathcal L$ fuera de $K$ tienen norma mayor que
$\norm{\theta_0}_2$, de modo que no pueden mejorar el mínimo alcanzado en $K$.
\end{block}
\end{frame}

\begin{frame}{Unicidad de norma mínima}
Supongamos que $u,v\in\mathcal L$ tienen la misma norma mínima $\mu$.
Como $\mathcal L$ es afín, $(u+v)/2\in\mathcal L$.
La identidad del paralelogramo da
\[
\left\lVert\frac{u+v}{2}\right\rVert_2^2
=\frac{\norm{u}_2^2+\norm{v}_2^2}{2}
-\frac{\norm{u-v}_2^2}{4}
=\mu^2-\frac{\norm{u-v}_2^2}{4}.
\]
Si $u\neq v$, el punto medio tendría norma menor que $\mu$, una contradicción.
Por tanto $u=v$.
\end{frame}

\begin{frame}{Coeficientes no únicos}
Considere dos columnas idénticas:
\[
A=\mat{1&1\\1&1\\1&1\\1&1},\qquad b=\mat{1\\2\\2\\4}.
\]
El ajuste depende únicamente de $\theta_1+\theta_2$. Su valor óptimo es
\[
\theta_1+\theta_2=\bar y=2.25.
\]
Entre todas estas soluciones, la de menor norma euclídea es
\[
\theta_1=\theta_2=1.125.
\]
La independencia de las columnas determina si podemos identificar los coeficientes.
\end{frame}

\begin{frame}{Ajuste polinómico}
Para una parábola $\widehat y=c+ax+dx^2$,
\[
A=\mat{1&x_1&x_1^2\\\vdots&\vdots&\vdots\\1&x_m&x_m^2},
\qquad\theta=\mat{c\\a\\d}.
\]
El problema conserva la misma forma:
\[
\min_\theta\norm{A\theta-b}_2^2.
\]
Agregar columnas puede reducir el error de ajuste dentro de la muestra.
La capacidad de predecir fuera de ella requiere una evaluación adicional.
\end{frame}

\begin{frame}{Actividad de clase}
Para los datos
\[
(x_i,y_i)=(0,1),(1,2),(2,2),(3,4),
\]
\begin{enumerate}
\item Construya el sistema para ajustar $c+ax+dx^2$.
\item Calcule $A^TA$ y $A^Tb$.
\item Justifique la unicidad sin calcular una inversa.
\item Explique por qué el error mínimo no puede superar $0.70$.
\item Compare el ajuste cuadrático con la recta en $x=4$.
\end{enumerate}
\end{frame}

\begin{frame}{Cierre de la clase}
\begin{itemize}
\item La matriz de diseño representa la familia de funciones elegida.
\item Las ecuaciones normales caracterizan todos los minimizadores.
\item El residuo óptimo es ortogonal al espacio columna.
\item El rango completo garantiza coeficientes únicos.
\end{itemize}
\vspace{.4cm}
En la siguiente clase resolveremos el problema mediante QR y estudiaremos
el efecto del condicionamiento.
\end{frame}
\end{document}
